\subsection{E.2. Effectivity, slope cones, and the finiteness axiom}

\paragraph{E.2a. The slope description (Proved for Weil characters;
proved for numerical motives under the Tate conjecture).}
\label{E2a}
Let $L/\mathbb Q$ be a finite Galois CM field and let $M=M(q,L)$
consist of $p$-units $\lambda\in L^\times$ with
$\lambda\bar\lambda=q^{w(\lambda)}$, $w(\lambda)\in\mathbb Z$.
For a prime $\mathfrak p$ over $p$ put
\[
 i_{\mathfrak p}(\lambda)=
 \frac{\operatorname{ord}_{\mathfrak p}(\lambda)}
      {\operatorname{ord}_{\mathfrak p}(q)}\in\mathbb Q.
\]
These are rational cocharacters of $D(M)^0$, satisfying
$i_{\mathfrak p}+i_{\tau\mathfrak p}=w$ and
$i_{\mathfrak p}(q)=1$. With $M^+=M\cap\mathcal O_L$,
\[
 M^+=\bigcap_{\mathfrak p\mid p}
       \{\lambda\in M:i_{\mathfrak p}(\lambda)\geq0\}.
\]
Thus arithmetic effectivity is the intersection of the halfspaces
of the Galois conjugate rational slope cocharacter. For a rational
representation, the condition is imposed on all its geometric
characters. At that level this is the assertion denoted by the
intersection of the $P(\sigma i)$ in the question; $i$ must be
allowed to be rational.

\begin{proof}
An element of $M$ has valuation zero at every finite prime outside
$p$. It is an algebraic integer exactly when its remaining
valuations are nonnegative. The two cocharacter identities follow
by taking valuations in $\lambda\bar\lambda=q^{w(\lambda)}$.
The denominators are harmless for a sign condition and can be
cleared simultaneously on a finite-rank character lattice.

For actual numerical motives we make the extra hypothesis explicit.
Assume the Tate conjecture in the form giving the classification of
simple numerical motives by conjugacy classes of Weil numbers
\citep[Proposition~2.6]{Milne1994}. An effective motive has integral
Frobenius eigenvalues by the Weil theorem. Conversely, decompose a
motive with integral eigenvalues into simple numerical motives. For
each eigenvalue B.1a supplies an effective abelian or Artin motive
with that eigenvalue. A simple constituent with the chosen
conjugacy class is isomorphic to the original simple motive by the
assumed classification. It is effective, because direct factors of
effective numerical motives are effective. Taking the direct sum
proves the converse. We do not drop this Tate hypothesis: equality
of eigenvalues alone does not give an isomorphism of arbitrary
motives without the required full-faithfulness statement.
\end{proof}

\paragraph{E.2b. A single normalized integral cocharacter for all motives
(Refuted for every $q$).}
\label{E2b}
For the full Weil-number group $M(q)$ there is no homomorphism
$i_1:M(q)\to\mathbb Z$ with $i_1(q)=1$. In particular the global
arithmetic effectivity structure cannot be $P(i_1)$ for a class in
the specified $I(G)$ with $\varepsilon i_1=\mathrm{id}$.

\begin{proof}
The number $a=\sqrt q$ is an algebraic integer all of whose conjugates
have absolute value $q^{1/2}$ and is a $p$-unit. Thus $a\in M(q)$
and $q=a^2$. A group homomorphism to $\mathbb Z$ would give
$1=i_1(q)=2i_1(a)$, an impossibility. Honda--Tate realizes $a$ on
$H^1$ for every $p$ and $q$, so this is a counterexample in the
motivic range, not an extraneous formal character. If $q$ is a
square the Honda--Tate class is supersingular elliptic; if $q$ is
not a square it has dimension two by
\citet[Th\'eor\`eme~1 and examples, pp.~96--97]{Tate1971}.
\end{proof}

\paragraph{E.2c. Finite subcategories and an integral criterion (Proved).}
\label{E2c}
For a Galois-stable finitely generated character group $N\subset
M(q,L)$ containing $q$, the distinguished slope at $\mathfrak p$
is integral exactly when
\[
 \operatorname{ord}_{\mathfrak p}(N)
       \subset\operatorname{ord}_{\mathfrak p}(q)\mathbb Z.
\]
This follows immediately from its definition as a homomorphism
$N\to\mathbb Q$; its values must lie in $\mathbb Z$ to be a
cocharacter. It is a criterion for this specified slope, not a
claim that all normalized integral cocharacters have been classified.
The example from R8 is confirmed exactly:
\[
 L=\mathbb Q(i),\quad q=25,\quad
 \lambda=15+20i=(2+i)^3(2-i).
\]
At the two primes above $5$ its normalized slopes are $3/2$ and
$1/2$ although its weight is $2$. Thus even weights alone do not
make the slope integral. The supplied program verifies the
factorization and norm with exact Gaussian arithmetic.

\paragraph{E.2d. Recovery of the rational slope orbit (Proved;
precise form of Lead E1).}
\label{E2d}
For the full finite-field lattice $M(q,L)$, its arithmetic effective
cone together with $q$ recovers the Galois orbit of the normalized
rational slope cocharacter. If a prime over $p$ splits in $L/K$,
the stabilizer of a recovered slope is its decomposition group and
its field of definition is $L^{D_{\mathfrak p}}$, determined up to
$\mathbb Q$-conjugacy. No integral cocharacter is asserted.

\begin{proof}
Since $L/\mathbb Q$ is Galois and complex conjugation is central,
either all primes over $p$ are fixed by conjugation, or none are.
In the latter case choose one prime in each of $s$ pairs. On
$M/\mu(L)\otimes\mathbb R$ use coordinates
$(w,a_1,\ldots,a_s)$, where $a_j=i_{\mathfrak p_j}$.
They are independent: B.4a gives a full lattice of anti-divisors
after multiplication by $2h$, and the character $q$ supplies the
weight direction. The effective cone is therefore exactly
\[
 \{(w,a_1,\ldots,a_s):0\leq a_j\leq w\quad(1\leq j\leq s)\}.
\]
Its $2s$ facets have inward normal rays $a_j$ and $w-a_j$.
Each is a genuine facet: keep its equality and make all other
inequalities strict. These are exactly the primes-over-$p$ slopes.
Normalizing every normal by its value $1$ on $q$ removes the positive
scaling ambiguity. Galois acts on the cone and its facets and acts
transitively on the set of primes. Its stabilizer of a specified
normal is exactly $D_{\mathfrak p}$, since the normals are distinct.
This recovers their orbit and its field of definition. If every
prime is fixed by conjugation, the norm-one lattice has rank zero,
all slopes equal $w/2$, and the recovered orbit is the singleton
defined over $\mathbb Q$.
\end{proof}

\paragraph{E.2e. The finiteness axiom (Refuted globally; finite-lattice
version proved).}
\label{E2e}
The marginal proposal on 14:3, ``un nb fini de types de motifs
effectifs de poids $i$'', is false for the full category over every
$\mathbb F_q$, even when ``types'' means simple isomorphism classes.
Already in weight zero, for every $r\geq1$ the cyclotomic direct
summand of $h^0(\operatorname{Spec}\mathbb F_{q^r})$ has Frobenius
minimal polynomial $\Phi_r(T)$. The rational cyclic group algebra
is semisimple and its cyclotomic factor $\mathbb Q(\zeta_r)$ is a
simple rational representation, so these give pairwise nonisomorphic
simple effective Artin motives. This proof involves no conjecture
on cycles.

There are also infinitely many simple effective weight-one types.
The numbers $\sqrt q\,\zeta_r$ are Weil $q$-integers, and their
degrees are unbounded: the field generated by one of them together
with $\sqrt q$ contains $\zeta_r$, hence its degree is at least
$\varphi(r)/2$. Honda--Tate assigns distinct simple isogeny classes
along a subsequence of unbounded degrees. The same construction in
each positive weight, using B.1a and the spectral factor, gives
infinitely many eigenvalue types. The simple-motive classification
in all weights is not needed to refute the axiom.

For a fixed finite Galois CM field $L$ and a fixed weight $i$, however,
there are only finitely many effective characters in $M(q,L)$:
each valuation lies between $0$ and
$i\operatorname{ord}_{\mathfrak p}(q)$ and there are finitely many
prime valuations. A specified divisor and weight determine the
element up to an ordinary norm-one unit, which is a root of unity
of the fixed finite group $\mu(L)$. This proves the sharp
finite-lattice statement behind R13 while explaining why it cannot
be passed to the full protorus.
